Eine ideale Kühlmaschine entzieht dem kalten Raum (Temperatur $T_\text{k}$) die Wärme $Q_\text{k}$ und gibt mit der zugeführten Arbeit $W$ die Wärme $Q_\text{w} = Q_\text{k}+W$ an die wärmere Umgebung (Temperatur $T_\text{w}$) ab. Ihre Leistungsziffer ist (FoTaBe S.~172)
\begin{align}
 \eta_\text{K} &= \frac{Q_\text{k}}{W} = \frac{T_\text{k}}{T_\text{w}-T_\text{k}}
\end{align}
Sie ist grösser als eins, wenn
\begin{align}
 \frac{T_\text{k}}{T_\text{w}-T_\text{k}} &> 1 \\
 T_\text{k} &> T_\text{w}-T_\text{k} \\
 T_\text{k} &> \frac{1}{2}\,T_\text{w}
\end{align}
Die absolute Temperatur des kalten Raums muss also mehr als halb so gross sein wie die absolute Temperatur der Umgebung. Bei Zimmertemperatur liegt diese Grenze weit unter dem Gefrierpunkt; jeder gewöhnliche Kühlschrank und jede Tiefkühltruhe erfüllt die Bedingung.
